\documentclass[11pt,a4paper]{article}
\usepackage[margin=25mm]{geometry}
\usepackage{amsmath,amssymb,amsthm,booktabs}
\newtheorem{lemma}{Lemma}
\newcommand{\Aut}{\operatorname{Aut}}
\newcommand{\C}{\operatorname{C}}
\title{Outer reflexibility when one rotation has order three}
\author{Human Compression project}\date{}
\begin{document}\maketitle
Let $G=O'N$ and $A=\Aut(G)=G{:}2$. We count $A$-orbits of ordered
pairs $(R,S)$ satisfying
\[
 |R|=3,\qquad |S|=q>2,\qquad |RS|=2,\qquad\langle R,S\rangle=G,
\]
whose simultaneous inverter is outer. Source equation (1.1) on printed
page 212 prescribes these generating pairs, and Theorem 1.1(2) on
page 213 takes their orbits under $\Aut(G)$. We use the same convention
as the previously established generating-pair and inversion arguments;
there is no division by duality. Since $3\ne15$, choosing the order-three
rotation first also identifies the source's unordered notation uniquely
in the case corrected below. We give a finite certificate for every such $q$.
In particular, the count for $(3,15)$ is one, whereas the fixed source
states two. This does not decide its separate assertions about inner
inverters or chirality for that type.

\section{Two elementary reductions}
The acquired character calculation supplies one class of outer involutions,
a representative $a$, and
\[
 |\C_A(a)|=351120,\quad
 \kappa_3=|\{b\in a^A:|ab|=3\}|=5852,\quad
 \kappa_2=|\{c\in a^A:|ac|=2\}|=1463.
\]
Put $J=\C_G(a)$. Since $a$ lies outside the index-two subgroup $G$,
\[
 \C_A(a)=\langle a\rangle\times J,\qquad |J|=175560.
\]
No identification of $J$ as an abstract named group is needed below.

\begin{lemma}\label{colors}
Suppose $J_0\le J$ is given by words in the acquired generators and acts
on a finite set $X$. Let $\Omega$ be a $J_0$-orbit in $X$, and label the
$J_0$-orbits in $X$ by distinct colors. For an outer involution $b$ with
$|ab|=3$, color $x\in\Omega$ by the orbit containing $x^b$.
Suppose this coloring has an orbit of size $5852$ under the image of
$J_0$ on $\Omega$, and its stabilizer contains an explicitly verified
subgroup of order $30$. Then $J_0=J$, the action on $\Omega$ is faithful,
and the coloring stabilizer has order $30$.
If lifts of its displayed generators centralize $b$, then
$H=\C_J(b)$ has order $30$ and $b^J$ is the entire set counted by
$\kappa_3$.
\end{lemma}
\begin{proof}
Write $\overline J_0$ for the image on $\Omega$. Orbit--stabilizer gives
$|\overline J_0|\ge5852\cdot30=175560$. On the other hand,
$|\overline J_0|\le |J_0|\le |J|=175560$. Equality holds throughout.
It proves the first conclusions and also that the displayed subgroup is
the full coloring stabilizer.

An element of $J$ commuting with $b$ preserves this coloring, because it
preserves each $J$-orbit in $X$. Conversely the specified generators of
the full coloring stabilizer have lifts commuting with $b$. Faithfulness
therefore identifies that stabilizer with $\C_J(b)$. Thus
$|b^J|=175560/30=5852$. This orbit is contained in the set of outer
involutions whose product with $a$ has order three, which has exactly
$5852$ elements. It is consequently the complete set.
\end{proof}

\begin{lemma}\label{generation}
Let $G$ act on $X$, let $\Delta$ be a $G$-orbit of size $n$, and suppose
no proper maximal subgroup of $G$ has order divisible by $n$.
For $K\le G$ and $x\in\Delta$, one has $K=G$ if and only if
$|x^K|=n$.
\end{lemma}
\begin{proof}
The forward implication is immediate. In the reverse direction,
orbit--stabilizer gives $n\mid |K|$. If $K<G$, choose a proper maximal
subgroup containing $K$. Its order would be divisible by $n$, contrary
to the hypothesis.
\end{proof}

\section{The acquired finite certificate}
The named faithful permutation model of $A$ has degree $245520$.
Its standard generators $a,d$ have orders $2,4$, with $|ad|=22$;
$a$ is outer and $d,d^a$ generate the socle. Those identifications are
borrowed premises of the already accepted model. The literal reader uses
the exact same permutation files and checks the displayed orders.

Eight retained words give elements of $J$. They were obtained using the
usual Bray expression $h(aa^h)^{(k-1)/2}$ when $|aa^h|=k$ is odd.
The certificate contains the full words. A repeated-subword factorization
speeds their evaluation, and its equality with each full word is checked
as a literal list before evaluating any permutation. Their commutation
with $a$ and their membership in the socle follow respectively from
literal multiplication and the even number of outer-generator letters.
The product-order-three partner is also an explicit conjugate $b=a^h$,
with $h\in A$ given by a retained word.

The literal permutation reader finds the following orbit sizes for the
eight generators on $X$:
\[
 1540^2,\quad4180^2,\quad14630^6,\quad29260^2,\quad43890^2.
\]
Exponents denote multiplicities. It uses one complete orbit $\Omega$ of
size $1540$. It checks every generator image on that orbit, constructs
the coloring from the actual full-domain orbits, and enumerates exactly
$5852$ distinct colorings by the small generator permutations. Two
explicit permutations generate a group of order $30$ preserving the
coloring. Their complete closure is enumerated, and their lifted words
commute with both $a$ and $b$. Thus Lemma~\ref{colors} applies without a
new group-order or centralizer algorithm as a premise.

The reader also constructs a conjugacy orbit of $1463$ involutions in
this faithful image of $J$. This is the entire involution set: the map
$j\mapsto aj$ identifies involutions of $J$ with the outer involutions
$c\ne a$ commuting with $a$, whose number is $\kappa_2=1463$.
The literal $H$-action partitions this set into $55$ disjoint orbits.
Each representative has a word in the eight generators, and its lift is
checked to have order two, commute with $a$, and restrict to that precise
small permutation. No conjugacy exhaustion returned by GAP is needed
for these final checks.

For generation, the reader computes the orbit of one point under the
socle generators $d,d^a$. Its size is $122760$, not the full degree of
the $A$-action. The complete acquired list of maximal-subgroup orders is
\[
 3753792^2,\quad175560,\quad161280,\quad25920^2,\quad14880^2,
 \quad10752,\quad7920^2,\quad2520^2.
\]
None is divisible by $122760$. Lemma~\ref{generation} therefore decides
whether each pair generates $G$ by a literal orbit calculation. The
reader enumerates the reachable points and checks invariance under both
generators. Completeness of the maximal list remains an explicit
borrowed subgroup-classification premise; it is not inferred from the
absence of additional orders in a computation.

\section{Counting the outer pairs}
For an involution $j\in J$ set $c=aj$, $R=ab$, and $S=bc$.
Then $a,b,c$ are outer involutions, $|R|=3$, and $RS=ac=j$ has
order two. The accepted inversion correspondence gives a bijection
between the generating triples and the generating pairs with an outer
simultaneous inverter. It also proves that their $A$-action is free.

By Lemma~\ref{colors}, all eligible $b$ form one $J$-orbit, hence one
$\C_A(a)$-orbit. The stabilizer $T=\C_A(a,b)$ has order
$351120/5852=60$. Thus, after fixing $a,b$, the desired orbit count is
\[
 \frac1{60}\sum_{\mathcal O}|\mathcal O|,
 \tag{1}
\]
where the sum runs over the $H$-orbits whose representative satisfies
$|bc|=q$ and $\langle ab,bc\rangle=G$. These properties are constant
on each $H$-orbit. Formula (1) does not require explicitly fusing those
orbits under $T$: their union is $T$-invariant and $T$ acts freely on
it. In particular the denominator is $60$, not $30$.

\begin{center}\begin{tabular}{r r r r}\toprule
$q$&All involutions $j$&Generating $j$&Outer pair orbits\\\midrule
10&150&60&1\\
11&150&120&2\\
14&210&180&3\\
15&120&60&1\\
19&300&180&3\\
20&120&120&2\\
28&120&120&2\\
31&60&60&1\\
\bottomrule\end{tabular}\end{center}


The auxiliary values $q=2,5,6,7$ have no generating representatives.
There are no representatives of orders $3,4,8,12,16$ in the complete
involution partition. Together with the table, this accounts for every
rotation order in the fixed source. The total number of outer classes
with first order three is $15$.

For $(3,15)$, the four $H$-orbits each have size $30$. Their point-orbit
sizes are respectively $14630,122760,122760,43890$ (certificate indices
$2,23,38,51$). Exactly two generate $G$, so (1) gives $60/60=1$.
This corrects the statement in the proof of source Lemma 5.1 on printed
page 218 (physical page 8), which assigns two of the six generating-pair
orbits to outer inverters. The numerical total six is unchanged.
The source's concluding totals imply $284-230-39=15$ outer classes in
this row, which agrees with our total; this agreement does not establish
its separate chiral or inner totals.

The inversion correspondence and earlier character counts are reused
proved inputs. Orbit--stabilizer, conjugation, the Bray construction and
the maximal-subgroup criterion are standard arguments; no priority claim
is made. The new contribution here is their application with literal
color, orbit and generation certificates to this complete source row.
The other rows, the inner spectrum, chirality, and whole-source completion
remain separate obligations. No conjecture or Lean credit is asserted.
\end{document}
